\begin{table}[htbp]
\centering
\caption{Sensitivity to covariate imputation}
\label{tab:robustez_sem_imputacao}
\resizebox{\textwidth}{!}{%
\begin{tabular}{lccccc}
\toprule
Grade-subject & DR-DiD (main) & DR-DiD no imputation & Main N & No-imputation N & N fraction \\
\midrule
9th LP & 0.510 & 0.384 & 24754 & 10804 & 0.436 \\
9th Mat & 0.601 & 0.502 & 24769 & 10819 & 0.437 \\
HS3 LP & 0.571 & -- & 27515 & -- & -- \\
HS3 Mat & 0.627 & -- & 27530 & -- & -- \\
\bottomrule
\end{tabular}
}
\vspace{0.2cm}
\begin{flushleft}
\footnotesize Note: the no-imputation version uses only observations with complete raw covariates. The exercise should be read as a diagnostic of sensitivity and sample coverage, not as a direct substitute for the main specification.
\end{flushleft}
\end{table}



